%%%
%
% $Autor: c. Crety, J. Perewersenko, H.Mey $
% $Datum: 2025-06-14 $
% $Pfad: ArduinoIDE2 $
% $Dateiname: 
% $Version: 4620 $
%
% !TeX spellcheck = de_DE/GB
% !TeX program = pdflatex
% !BIB program = biber/bibtex
% !TeX encoding = utf8
%
%%%


\chapter{Arduino IDE 2.3.6}\label{ArduinoIDE2}



In diesem Kapitel werden die Funktionen und die Benutzeroberfläche der Arduino IDE erklärt. Die einzelnen Komponenten werden erklärt und es wird aufgezeigt, wie man diese benutzt. Des Weiteren gibt es eine Installationsanleitung und es wird aufgezeigt, wie der Arduino Nano 33 BLE Sense mit der Arduino IDE verknüpft wird.


\section{Beschreibung der Arduino IDE}\label{ArduinoIDE}

Die Arduino IDE ist eine open-source Software zur Bearbeitung, Kompilierung und Übertragung von Codes auf Arduino-Module. Sie ist plattformübergreifend nutzbar, basiert auf Java und ist mit Windows, macOS und Linux kompatibel. Die Arduino IDE unterstützt verschiedene Arduino-Module, sowie die Programmierung in C und C++. 

Die Programme, welche in der Arduino IDE geschrieben werden nennt man "Sketches" und können auf den Mikrocontroller übertragen werden.


\bigskip

Wenn man die Arduino IDE öffnet, dann sieht man in der oberen Menüleiste \autoref{fig::Menu Button} 5 Schaltflächen:

\begin{center}

    \includegraphics[width=\textwidth]{Arduino/ArdiunoIDE2/menuboard}

    \captionof{figure} {Menu Button}\label{fig::Menu Button}		
\end{center}

Die Schaltflächen haben folgende Funktionen (von links nach rechts):

\begin{itemize}
	\item Verify: der Haken überprüft den Code auf Fehler und zeigt, ob er korrekt strukturiert ist.
	\item Upload: der Pfeil kompiliert den Code und überträgt ihn an das angeschlossene Board.
	\item Start Debugging: startet eine Debugging-Sitzung, um den Code Schritt für Schritt durchzugehen.
	\item Serial Plotter: Echtzeit-Daten vom Board werden grafisch angezeigt.
	\item Serial Monitor: Terminalfenster zum Senden und Empfangen von Textdaten in Echtzeit über die serielle Schnittstelle.
\end{itemize}


\section{Installation of the Arduino IDE}

\subsection{Download of the Arduino IDE}

Der Arduino Nano 33 BLE Sense verwendet die  Arduino Softwareumgebung \textbf{Integrated Development Environment (IDE)} zur Programmierung. Diese IDE ist die am weitesten verbreitete für alle Arduino-Boards und kann sowohl online als auch offline verwendet werden. Die Arduino IDE ist eine open-source Software und vereinfacht das Schreiben von Code und das Hochladen auf das Board erheblich. Die Software ist für Windows, Linus und macOS verfügbar und die aktuellste Version kann von der offiziellen Arduino-Website heruntergeladen werden: \HREF{https://www.arduino.cc/en/software}{Arduino-Software}

\section{Übersicht}
Die Arduino IDE 2 verfügt über eine Seitenleise, die den einfachen Zugriff auf die am häufigsten verwendeten Werkzeuge ermöglicht. ~\ref{ide2Overview} \cite{arduinodescription:2024}

\begin{itemize}

    \item \textbf{Verify / Upload:} Diese Funktionen dienen zum Kompilieren und Hochladen von Code auf ein Arduino-Board.
    \item \textbf{Select Board and Port:} Erkannte Arduino-Boards werden automatisch angezeigt, inklusive der zugewiesenen Portnummer.
    \item \textbf{Sketchbook:} Das Sketchbook ist ein zentraler Ort zur Verwaltung aller lokal gespeicherten Sketches. Es bietet eine übersichtliche Struktur zum Organisieren und Bearbeiten von Projekten und eine Synchronisierungsfunktion mit der Arduino Cloud, wodurch Sketches geräteübergreifend verfügbar sind.
    \item \textbf{Boards Manager:} Hier können offizielle Boardpakete von Drittanbietern installiert werden.
    \item \textbf{Library Manager:} Mit dem Bibliotheksverwalter hat man Zugriff auf tausente Arduino-Bibliotheken, welche von Arduino selbst oder der Community erstellt wurden.
    \item \textbf{Debugger:} Der Debugger ermöglicht das Testen und Debuggen von Programmen in Echtzeit.
    \item \textbf{Search:} Hier kann der geschriebene Coder auf Schlüsselwörter durchsucht werden.
    \item \textbf{Open Serial Monitor:} Hier wird der serielle Monitor als neuer Tab im unteren Konsolenbereich geöffnet.
\end{itemize}


\begin{center}
        \includegraphics[width=\linewidth]{Arduino/ArdiunoIDE2/ide2Overview.png}
        \captionof{figure}{Overview}\label{ide2Overview}
\end{center}


\section{Funktionen}
Die Hauptfunktionen der Arduino IDE sind die direkte Installation von Bibliotheken, die Synchronisierung von Sketches über die Cloud und die integrierten Debugging-Tools. Im Folgenden werden einige der zentralen Funktionen näher beschrieben.

\subsection{Sketchbook}

Das Sketchbook ~\ref{local-sketchbook} ist der Speicherort der Arduino-Code-Dateien, die als Sketches bezeichnet werden. Diese Dateien werden mit der Endung ".ino" gespeichert. Dabei ist wichtig, dass jeder Sketch in einem Ordner gespeichert werden muss, der genau denselben Namen hat wie die Sketch-Datei. Als Beispiel: Der Sketch heißt "MySketch.ino", also muss der Ordner, in welchem sich der Sketch befindet, "MySketch" benannt werden.

Standardmäßig werden Sketches in einem Ordner namens Arduino gespeichert, der sich im Dokumenten-Verzeichnis des Betriebssystems befindet.
Im linken Seitenmenü der Arduino IDE befindet sich ein Ordner-Symbol, über das man direkt zum Sketchbook gelangt. Dort können Sketches geöffnet, bearbeitet und verwaltet werden.

\begin{center}
        \includegraphics[width=0.7\linewidth]{Arduino/ArdiunoIDE2/localSketchbook.png}
        \captionof{figure}{Sketchbook}\label{local-sketchbook}
\end{center}


\subsection{Boards Manager}
Der Boards Manager ist über folgenden Pfad in der IDE erreichbar: \menu{Tools > Board > Board Manager}.

In der Arduino IDE ermöglicht der Boards Manager die Installation verschiedener Board-Pakete, wie in Abbildung ~\ref{boardManager} gezeigt.
Ein Board-Paket enthält alle notwendigen Dateien und Anweisungen, um Code für bestimmte Arduino-Boards zu kompilieren und hochzuladen.

Es sind verschiedene Board-Pakete verfügbar. Jedes Paket unterstützt verschiedene Arduino Boards, der Boards Manager stellt also sicher, dass die richtigen Tools und Compiler installiert sind und die Programmierung für das gewählte Board reibungslos funktioniert.

\begin{center}
        \includegraphics[width=0.7\linewidth]{Arduino/ArdiunoIDE2/boardManager.png}
        \captionof{figure}{Board Manager}\label{boardManager}
\end{center}


\subsection{Library Manager}
Der Library Manager ist über folgenden Pfad in der IDE erreichbar: \menu{Tools > Manage Libraries}.

In der Arduino IDE ermöglicht der Library Manager das Durchsuchen und Installieren einer Vielzahl von Bibliotheken. Bibliotheken erweitern die grundlegenden Arduino-Funktionen und erleichtern komplexe Aufgaben wie das Auslesen spezieller Sensoren oder die Ansteuerung von Modulen.
Diese Bibliotheken bieten vorgefertigten Code, eine Vereinfachung der Programmierung, die Unterstützung spezifischer Hardwarekomponenten und eine schnellere und effizientere Entwicklung. Der Library Manager ist in Abbildung ~\ref{librarymanager} zu sehen.

\begin{center}
        \includegraphics[width=0.7\linewidth]{Arduino/ArdiunoIDE2/libraryManager.png}
        \captionof{figure}{Library Manager}\label{librarymanager}
\end{center}

{

\section{Beispiele}

Ein zentraler Bestandteil der Arduino-Dokumentation ist die Vielzahl an Beispielsketches, die mit unterschiedlichen Bibliotheken mitgeliefert werden. Diese Beispiele demonstrieren die praktische Anwendung von Bibliotheksfunktionen und zeigen deren beabsichtige Verwendung und deren Hauptfunktionen im Einsatz. Die Beispielsketches können aus vorinstallierten Bibliotheken stammen oder aus den Board-spezifischen Bibliothekspaketen.
Zugriff auf die Beispiele erhält man über \menu{File > Examples}, und dann in der List die gewünscht Bibliothek auswählen.

\begin{center}
        \includegraphics[width=0.7\linewidth]{Arduino/ArdiunoIDE2/examplesketches.png}
        \captionof{figure}{Examples}\label{examplesketches}
\end{center}


In Abbildung  ~\ref{examplesketches} sieht man die Beispielliste, wenn ein Arduino UNO R4 WiFi an den Computer angeschlossen ist.



\section{Debugging}

Das Debugger Tool in der Arduino IDE dient dazu, Programme gezielt zu testen und zu debuggen. Zum Einen kann man mit dem Debugger Breakpoints setzen, um den Programmablauf an bestimmten Stellen zu unterbrechen. Zum Anderen gibt es die Funktion des schrittweisen Durchlaufend, um den Code Zeile für Zeile abzuarbeiten. Außerdem können Variablen überwacht werden, somit kann man Variablenwerte in Echtzeit einsehen, um Fehler besser einzugrenzen.
Das Debugging bietet Vorteile, wie das Auffinden von Logikfehlern, das Überprüfen von Zustandsänderungen und es bietet Einsicht in den Speicherzustand während der Programmausführung.
